\section{Глава~\ref{sec:chemoutput}. Химический выход}

Два провода к сердцу и их разная скорость, отрицательная обратная связь в гормональных каскадах, код частотой порций, суточный генератор и его подстройка светом измерены прямо; граница $K<8$ --- теорема теории регулирования, выведенная в главе; пульсации, порождаемые самой обратной связью каскада, --- гипотеза с сильным опытом в её пользу.

{\footnotesize
\setlength{\tabcolsep}{3pt}\renewcommand{\arraystretch}{1.15}
\begin{longtable}{>{\raggedright}p{0.5cm}>{\raggedright}p{4.0cm}>{\raggedright}p{5.6cm}>{\raggedright}p{2.3cm}>{\raggedright\arraybackslash}p{1.7cm}}
\toprule
№ & Утверждение & Опыт и что измерено & Источник & Статус \\
\midrule\endfirsthead
\toprule
№ & Утверждение & Опыт и что измерено & Источник & Статус \\
\midrule\endhead
1 & Ритмоводитель сердца сам бьётся около $100$ раз в минуту; в покое тормоз нажат, и частота обычно порядка $60$--$70$ & Люди, полное выключение обоих проводов к сердцу: собственная частота выше частоты покоя и убывает с возрастом, у взрослых порядка $100$ ударов в минуту. Значит, в покое преобладает тормоз. Частота покоя $60$--$70$ --- типичное значение, отдельно не цитировано. & \cite{JoseCollison1970ihr} & измерено \\
2 & Газ \nt{NORA} и тормоз \nt{ACET} --- фильтры нижних частот, тормоз быстрее~\eqref{eq:gasbrake} & Собака, частотно-модулированная стимуляция блуждающего или симпатического нерва сердца и передаточная функция к частоте предсердий: оба пути --- фильтры нижних частот, у симпатического частота среза много ниже и добавлена чистая задержка около $1.7\,\text{с}$. Характеристики зависят от среднего тонуса, так что линейная формула~\eqref{eq:gasbrake} --- приближение. & \cite{Berger1989} & измерено \\
3 & Тормоз быстр, потому что калиевый канал \nt{ACET} открывает без второго посредника, а \nt{NORA} --- через цепочку реакций & Мембранные лоскуты клеток предсердий: калиевый канал, открываемый ацетилхолином, открывают $\beta\gamma$-субъединицы G-белка, связанного с рецептором, без второго посредника внутри клетки. Путь \nt{NORA} через цАМФ здесь не цитирован. & \cite{Logothetis1987gbg} & измерено \\
4 & Кровь обходит тело за время порядка минуты; \nt{ADRE} из кровотока доходит до всех органов с рецептором & Общая оценка порядка величины; в главе так и сформулирована, источник не цитирован. & --- & оценка \\
5 & Гормональный каскад охвачен отрицательной обратной связью: конечный гормон тормозит первые ступени & Обзор опытов: гормоны коры надпочечников подавляют выброс АКТГ гипофизом и рилизинг-гормона гипоталамусом в трёх временных диапазонах --- быстром, промежуточном и медленном. & \cite{KellerWood1984fb} & измерено \\
6 & Три одинаковые ступени устойчивы при $K<8$~\eqref{eq:cascadestable}, на границе период $2\pi\tau/\sqrt3\approx3.6\tau$ & Выведено в главе: на частоте $\sqrt3/\tau$ три ступени дают сдвиг фазы $180^\circ$ и ослабление в $8$ раз. & вывод в главе & теория \\
7 & Ошибка уровня $1/(1+K)$, поэтому точнее примерно $10\,\%$ такой регулятор не держит (утверждение~\ref{prop:cascade}) & Следствие строки~6 для пропорциональной обратной связи. Настоящая ось имеет обратную связь на нескольких ступенях и в нескольких временных диапазонах (строка~5), так что граница относится к модели~\eqref{eq:cascade}, а не измерена. & вывод в главе & теория \\
8 & При $K=4$ и $7$ уровень устанавливается, при $K=12$ --- ровные пульсации с периодом $3.7$--$3.9\,\tau$ (рис.~\ref{fig:cascade}) & Расчёт по \texttt{models/\allowbreak chem\_\allowbreak models.py}: при $K=12$ последовательные периоды $3.9$, $3.7$, $3.8$, $3.7\,\tau$; при $K=4$ размах в конце счёта $0.003$. & \texttt{models/\allowbreak chem\_\allowbreak models.py} & модель \\
9 & Гормон стресса выходит пульсами около раза в час; пульсы порождает сама обратная связь, без отдельного генератора & Крысы, постоянное вливание рилизинг-гормона: АКТГ и кортикостерон колеблются с физиологической, почти часовой частотой --- ритмичный вход от гипоталамуса для пульсов не нужен. Модель гипофиз--надпочечник с запаздывающей обратной связью даёт такие колебания. & \cite{Walker2012glc, Walker2010} & гипотеза \\
10 & Получатель читает частоту порций, а не уровень & Обезьяны без собственного выброса рилизинг-гормона гонадотропинов: постоянное вливание не восстанавливает выброс гормонов гипофиза, порции раз в час восстанавливают; переход на постоянное вливание снова глушит ответ. Утомление рецептора как фильтр верхних частот --- толкование. & \cite{Belchetz1978} & измерено \\
11 & Разные частоты порций по-разному действуют на разные клетки одной железы & Те же обезьяны: учащение порций до $2$--$5$ в час снижает оба гормона гипофиза; урежение до одной за $3$~часа снижает один (ЛГ) и повышает другой (ФСГ), так что их отношение задаётся частотой. & \cite{Wildt1981gnrh} & измерено \\
12 & Суточный ритм порождает внутриклеточная отрицательная обратная связь с задержкой в несколько часов & Обзор: белки часовых генов с задержкой подавляют транскрипцию собственных генов; петля из транскрипции и трансляции даёт период около суток. & \cite{TakahashiJS2017} & измерено \\
13 & Главный генератор --- группа из десятков тысяч нейронов, синхронизирующая остальное тело & Обзор: надперекрёстное ядро --- главный суточный генератор; отдельные его нейроны в культуре --- самостоятельные генераторы, сеть их синхронизирует; у мыши около $10^4$ нейронов с каждой стороны. & \cite{Welsh2010scn} & измерено \\
14 & Собственный период у человека $24.18$~часа & Люди в тусклом свете при расписании, оторванном от суток: период ритмов мелатонина, температуры и кортизола в среднем $24.18$~часа у молодых и пожилых, с узким разбросом. & \cite{Czeisler1999} & измерено \\
15 & Свет подводит генератор как фазовая автоподстройка~\eqref{eq:pll}; нужен сдвиг около $11$~мин в сутки & Люди, одиночный яркий свет $6.7$~часа в разное время суток: кривая фазового ответа первого типа с размахом $5$~часов --- задержка до минимума температуры тела, опережение после. Уравнение~\eqref{eq:pll} --- стандартная модель; $11$~мин $=0.18$~часа --- арифметика. & \cite{Khalsa2003prc} & измерено \\
16 & Без света подстройки нет, и ритм уходит на собственный период & Полностью слепые люди без восприятия света, живущие в обычном обществе: примерно у половины ритмы мелатонина и кортизола идут со своим периодом, не совпадающим с сутками. & \cite{Sack1992blind} & измерено \\
\bottomrule
\end{longtable}}
